\chapter{Cp / Boundary Layer Tab}
\label{chap:cp}

The \menu{Cp / Boundary layer} tab provides a detailed,
\emph{XFoil}-style view of a \textbf{specific operating point}: one
airfoil, one Reynolds number and one angle of attack. It gathers in a
single window the pressure distribution (viscous \emph{and} inviscid)
and the airfoil geometry surrounded by its boundary layer.

\screenshot{20_tab_cp.png}{1.0}{%
    \menu{Cp / Boundary layer} tab for a NACA~2412 at
    $Re = 10^6$ and $\alpha = 6\textdegree$. Top: the pressure
    coefficient $C_p$ (upper surface in yellow, lower surface in cyan)
    and its inviscid value (white dashed line). Bottom: the airfoil to
    scale, with the boundary layer edges $\delta^*$ on the upper
    surface (yellow) and lower surface (cyan). The inset summarizes the
    coefficients at the selected point.}

\section{Selecting the operating point}

A control bar at the top of the tab targets the calculation to
display:

\begin{tabularx}{\linewidth}{l X}
    \toprule
    \textbf{List} & \textbf{Role} \\
    \midrule
    \menu{Airfoil}
        & Airfoil to display: current (blue), reference (red) or
          flap (green). Only the airfoils actually simulated are
          offered. \\
    \menu{Re}
        & Reynolds number of the calculation. \\
    \menu{$\alpha$}
        & Angle of attack (in degrees) of the calculation. \\
    \bottomrule
\end{tabularx}

The lists are \textbf{cascading}: changing the airfoil updates the
list of available Reynolds numbers, then that of the angles of attack.

\begin{noteinfo}
Unlike the \menu{Results} tab, which overlays the airfoils, this tab
displays \textbf{a single operating point at a time}, for a detailed
reading of the pressure and the boundary layer.
\end{noteinfo}

\section{Top plot: pressure coefficient}

The top plot follows the \emph{XFoil} convention: the vertical axis
carries the pressure coefficient $C_p$ \textbf{pointing downward} —
suction ($C_p < 0$) is therefore directed \emph{upward}. Three curves
are overlaid:

\begin{itemize}
    \item the \textbf{upper surface} (yellow);
    \item the \textbf{lower surface} (cyan);
    \item the \textbf{inviscid} $C_p$ (white dashed line), obtained
        from an XFoil calculation without boundary layer (panel
        method, Reynolds-independent).
\end{itemize}

The gap between the viscous and inviscid $C_p$ reveals the effect of
the boundary layer: small at low angles of attack, it deepens near the
trailing edge (recompression) and in case of separation.

\section{Bottom plot: airfoil and boundary layer}

The bottom plot shows the airfoil \textbf{to chord scale} (isometric
frame: the airfoil proportions are preserved). Around the surface (in
white) are drawn the \textbf{boundary layer edges} — the surface
offset by the displacement thickness $\delta^*$ along the outward
normal:

\begin{itemize}
    \item upper surface in yellow;
    \item lower surface in cyan.
\end{itemize}

The thickening of these edges toward the trailing edge reflects the
growth of the boundary layer; a marked bulge signals separation. The
wake, downstream of the trailing edge, is excluded from the plot.

\section{Coefficients inset}

An inset, at the top right, recalls the integral coefficients at the
selected point: Reynolds, angle of attack $\alpha$, lift $C_L$, moment
$C_M$, drag $C_D$ and lift-to-drag ratio $C_L/C_D$.

\begin{astuce}
Compare the viscous $C_p$ with the inviscid dashed line to assess at a
glance the influence of viscosity at the chosen angle of attack; then
inspect the $\delta^*$ edges to spot separation starting at the
trailing edge.
\end{astuce}

\begin{noteinfo}
The inviscid $C_p$ is computed by a separate XFoil pass (without
boundary layer). This pass is very fast and can be disabled via the
\code{INVISCID\_CP} parameter (enabled by default).
\end{noteinfo}
